\section{Results}\label{sec:mp_results}

\subsection{Early prediction and target correction}

On the feasible-deployment development block, observable exposure reduced the
eight-point holdout RMSE for the objective mean from 0.594 to 0.270 and for the
constraint mean from 0.204 to 0.141.  Constraint-ranking Spearman correlation
rose from 0.119 to 0.624.  This is a prediction result; the initial true
recommendation was feasible in 7/8 tasks for both representations and their
mean initial rewards were essentially identical.

The risk factor-family prior also improved total log-risk prediction.  In the
source leave-one-task-out check its mean log predictive density gain over a
pooled-risk control was 0.692, and all eight development targets had lower
initial $\eta$ RMSE (0.301 versus 0.363).  The coefficient RMSE, however, was
not lower (0.255 versus 0.251), so the evidence supports total-risk calibration
rather than separate physical-factor recovery.

\subsection{Decision attribution}

The observable representation versus fixed policy-coordinate closed loop gave
mean final true rewards 0.8613 and 0.8538, respectively.  Only 3/8 tasks had a
positive search increment, so the predeclared stable-deployment gate failed.
Adding the learned risk prior to observable exposure reduced mean final reward
slightly to 0.8608; one task declined and the others were unchanged.

The new same-source control isolates the mean structural prior on the
development block:
\begin{table}[t]
\centering
\caption{Paired source-structure attribution on the eight-task development block.}
\label{tab:mp_structure_control}
\begin{tabular}{lrrrr}
\toprule
Arm & Tasks & Mean final reward & Certified & False certificates\\
\midrule
Source-selected structure & 8 & 0.86127 & 8/8 & 0\\
Same-source full dictionary & 8 & 0.85582 & 8/8 & 0\\
\bottomrule
\end{tabular}

\end{table}

The paired final-reward differences were positive on 3 tasks, negative on 1,
and zero on 4, with mean $+0.00545$.  Thus source structure has a measurable
mean-prior contribution in this block, but not a stable deployment advantage.

\subsection{Independent fixed-seed confirmation}

The separate v2 block retained all eight predeclared law seeds after its
oracle-only preflight.  The source-selected arm had mean final true reward
0.44274 (5/8 certified), while the same-source full dictionary had 0.45160
(4/8 certified); neither arm produced a false certificate.  The paired
differences were positive on 2 tasks, negative on 2, and tied on 4.  Thus the
development-block ordering is not confirmed by this independent block.  The
original 95433--95440 interval was not run as a block; seed 95434 overlaps the
development artifact and is excluded from any clean-confirmation claim.  The
v2 block is a separate contract and does not repair that overlap retroactively.

\begin{table}[t]
\centering
\caption{Independent fixed-seed confirmation block; all eight tasks retained.}
\begin{tabular}{lrrrr}
\toprule
Arm & Tasks & Mean final reward & Certified & False certificates\\
\midrule
Source-selected structure & 8 & 0.44274 & 5/8 & 0\\
Same-source full dictionary & 8 & 0.45160 & 4/8 & 0\\
\bottomrule
\end{tabular}
\end{table}
\par\smallskip
All eight fixed seeds were retained after the oracle-only preflight; paired
final-reward differences were positive on 2, negative on 2, and tied on 4 tasks.


\subsection{Same-contract target-only baseline}

The target-only full-dictionary control uses the same 68-function dictionary,
candidate bank, initial design, paid-observation budget, keyed noise streams,
pooled factor-family risk arm, acquisition rule, and independent verifier; it
starts without source labels.  Its mean final true rewards are 0.76199 on the
development block and 0.45148 on v2, compared with 0.86127 and 0.44274 for
source-selected structure.  The corresponding same-source full-dictionary
means are 0.85582 and 0.45160.  The source-selected minus target-only
difference is therefore +0.09928 in development but -0.00873 in v2.  This
feasible nonempty baseline completes the information-contract comparison and
does not restore an independent deployment claim.

We also ran a generic target-only RBF-KG baseline with fixed isotropic
lengthscale 0.25 and outputscale 1.0, using the same observable coordinates,
budget, keyed streams, and verifier but no source labels or source-risk prior.
Its means are 0.86239 on development and 0.45912 on v2, exceeding the
source-selected arm by -0.00112 and -0.01638, respectively.  Thus the source
structure is attributable relative to the deliberately full dictionary in the
development block, but it is not superior to a generic target-only RBF control.

\begin{table}[t]
\centering
\caption{Same-contract target-only baseline and paired mean differences.}
\label{tab:mp_target_only_main}
\begin{tabular}{lrrrrr}
\toprule
Block & Source-selected & Full dictionary & Target-only & RBF-KG & Source--RBF\\
 & mean (certified) & mean (certified) & mean (certified) & mean (certified) & mean difference\\
\midrule
Development & 0.86127 (8/8) & 0.85582 (8/8) & 0.76199 (7/8) & 0.86239 (8/8) & -0.00112 \\
Independent v2 & 0.44274 (5/8) & 0.45160 (4/8) & 0.45148 (4/8) & 0.45912 (4/8) & -0.01638 \\
\bottomrule
\end{tabular}

\end{table}

\subsection{Independent verification and negative evidence}

All 24 cells of the observable/risk development loop and all 16 cells of the
source-structure control received valid independent certification with no
false certificates.  This verifies the task and verifier contract; it does
not turn the development block into confirmation.  The original physical-risk
closed loop had seven of eight targets with no feasible candidate in the fixed
bank, and the official HyperBO comparison used a target with no feasible
candidate.  Those failures remain part of the result because candidate
coverage and task feasibility are part of an end-to-end deployment claim.

The risk-only A45 diagnostic found that refreshed mixture weights, moment
collapse, and quadrature-node choices changed probe rankings only modestly; it
did not identify the current Gaussian fantasy approximation as the cause of
the weak deployment increment.
